\clearpage
\VerdierPhysicalPageAnchor{0027}
$\operatorname{Ext}_{\mathcal A}^{n}(X,Y)$ の元、すなわち [1] の意味での
複体 $X,Y$ の第 $n$ 超 Ext と解釈できる。$X,Y$ が有界複体（有限個の
成分だけが非零）ならば、$X$ から $Y$ への次数 $n$ の準同型は Yoneda
構成 [3] を一般化する構成によって定義できる。$\mathbf D(\mathcal A)$
における射の合成は Yoneda 合成にほかならない。

\VerdierSemanticAnchor{intro-6-plan-chapter-three-spectral-objects}
第 4 節では通常のスペクトル対象を研究する。$\mathcal A$ の各複体は
$\mathbf D(\mathcal A)$ に値をとる三つのスペクトル対象を生み、それら
から主要なスペクトル系列が自動的に書き下される。さらに、複体と複体の射
の任意の列 $\cdots\to X_n\to X_{n+1}\to\cdots$（推移射は単射でなくても
よい）はスペクトル対象を生む。

\VerdierSemanticAnchor{intro-7}
\noindent\textbf{7.}\quad
要するに、われわれは一つの形式体系の基礎を築こうとしている。この種の
仕事には、省略三段論や連鎖三段論、数学的には真の困難を含まないものの
ときに（むしろ常に）退屈な確認の連鎖を要する証明など、あらゆる不都合が
伴う。したがって、この形式体系が柔軟性と一般性によって「真の」定理の
定式化と証明を可能にする限りにおいてのみ、本論文には意義がある。
コホモロジー双対定理を述べ、証明したいという要請が、この形式体系の最初
の動機であった。さまざまな状況において、形式的に非常によく似た
コホモロジー双対理論が存在する：

\VerdierSemanticAnchor{intro-7-duality-theories}
\begin{enumerate}
\renewcommand{\labelenumi}{\arabic{enumi})}
\VerdierSemanticAnchor{intro-7-duality-1}
\item スキーム上の連接層のコホモロジーにおける双対理論 [4]。

\VerdierSemanticAnchor{intro-7-duality-2}
\item スキームのエタール位相における層のコホモロジーの双対理論 [5], [6]。

これらの理論は A. Grothendieck による。

\VerdierSemanticAnchor{intro-7-duality-3}
\item J. Tate と J.-L. Verdier による Galois 加群のコホモロジーの
双対理論 [7]。

\VerdierSemanticAnchor{intro-7-duality-4}
\item 局所コンパクト空間上の層のコホモロジーの双対理論 [9], [10]。

\VerdierSemanticAnchor{intro-7-duality-5}
\item 連接解析層のコホモロジーの双対理論。この理論は現在も本質的に予想
の段階にある。

\VerdierSemanticAnchor{intro-7-duality-6}
\item 局所体、数体または関数体の双対理論と類体論との関係。この理論は
なお展開されるべきであり、その第一歩は M. Artin による [8]。
\end{enumerate}

\VerdierSemanticAnchor{intro-7-adjunction}
これらすべての理論は、導来圏の間のある関手を結ぶ随伴性を出発点
（または到達点）とする。
